\documentclass[11pt]{article}

% --- Standard packages ---
\usepackage{amsmath,amssymb}
\usepackage{graphicx}
\usepackage{booktabs}
\usepackage[colorlinks=true,linkcolor=blue,citecolor=blue,urlcolor=blue]{hyperref}
\usepackage{natbib}
\usepackage[ruled,vlined]{algorithm2e}
\usepackage{subcaption}
\usepackage{xcolor}
\usepackage{siunitx}
\usepackage{multirow}
\usepackage{bm}

% --- Custom commands ---
\newcommand{\todo}[1]{\textcolor{red}{\textbf{[TODO: #1]}}}
\newcommand{\Rey}{\mathrm{Re}}
\newcommand{\Reytau}{\Rey_\tau}
\newcommand{\bhat}{\hat{b}}
\newcommand{\Shat}{\hat{S}}
\newcommand{\Omhat}{\hat{\Omega}}
\newcommand{\nueff}{\nu_\mathrm{eff}}
\newcommand{\nusgs}{\nu_\mathrm{sgs}}
\newcommand{\nut}{\nu_t}

\title{Cost-Accuracy Tradeoffs for Neural Network Turbulence Closures\\in GPU-Accelerated Incompressible Flow Solvers}

\author{TBD}

\date{}

\begin{document}

\maketitle

\begin{abstract}
We systematically compare five data-driven turbulence closure architectures---multilayer perceptron (MLP), large MLP, tensor basis neural network (TBNN), physics-informed TBNN (PI-TBNN), and tensor basis random forest (TBRF)---against classical RANS models in a GPU-accelerated incompressible Navier--Stokes solver.
All models are trained on the same dataset \citep{McConkey2021} with the same case-holdout protocol and evaluated both \emph{a priori} (offline prediction accuracy) and \emph{a posteriori} (coupled CFD accuracy and wall-clock cost).
We provide a detailed computational anatomy of each closure, showing exactly which operations---matrix multiplies, activation functions, tree traversals, transport equation solves---dominate the cost.
On a \todo{GPU type} GPU, the TBNN achieves \todo{X}\% lower prediction error than $k$-$\omega$ SST at \todo{Y}\% additional computational cost per time step, while the smaller MLP adds only \todo{Z}\% overhead.
The TBRF yields the highest offline accuracy (validation RMSE of 0.064 vs.\ 0.085 for TBNN) but is impractical for GPU deployment due to branch-heavy tree traversals and poor cache utilization.
A key finding is that model rankings reverse between validation and test sets: the TBNN achieves the best validation accuracy but the worst test-set generalization (RMSE degrades $4.6\times$, with catastrophic failure on an unseen curved backward-facing step geometry), while the physics-informed PI-TBNN---previously thought to offer no benefit---degrades only $1.18\times$.
The realizability penalty acts as a regularizer that prevents overfitting to training geometries, yielding $2.7\times$ better test accuracy than the unconstrained TBNN at identical computational cost.
These results demonstrate that validation accuracy alone is a misleading metric for turbulence closure selection and provide practitioners with quantitative guidance for balancing predictive accuracy, generalization, and computational budget.
\end{abstract}

\section{Introduction}
\label{sec:introduction}

\subsection{RANS closure hierarchy}

Reynolds-averaged Navier--Stokes (RANS) simulations remain the workhorse of industrial computational fluid dynamics, enabling predictions of turbulent flows at a fraction of the cost of direct numerical simulation (DNS) or large-eddy simulation (LES).
The accuracy of RANS simulations hinges on the turbulence closure model, which approximates the Reynolds stress tensor $\overline{u'_i u'_j}$ that arises from the averaging process.
Closure models span a broad hierarchy of complexity and cost:
\begin{itemize}
    \item \textbf{Algebraic (zero-equation) models} compute the eddy viscosity from local mean-flow quantities---typically a mixing length and a velocity gradient---with $O(1)$ floating-point operations per cell.
    \item \textbf{One- and two-equation transport models} solve additional partial differential equations for turbulence quantities such as the turbulent kinetic energy $k$ and the specific dissipation rate $\omega$ \citep{Wilcox1988,Menter1994}. The $k$-$\omega$ SST model \citep{Menter1994} is the de facto standard in engineering practice, using blending functions to combine near-wall and free-stream formulations.
    \item \textbf{Explicit algebraic Reynolds stress models} (EARSM) augment transport models with tensor algebra to predict anisotropic Reynolds stresses without solving the full Reynolds stress transport equations.
    \item \textbf{Genetically evolved models} such as GEP \citep{Weatheritt2016} use symbolic regression to discover algebraic stress--strain relations from DNS data.
\end{itemize}
Moving up this hierarchy generally improves accuracy for complex flows (separation, secondary motions, streamline curvature) at the expense of additional computational cost and implementation complexity.

\subsection{Data-driven closures: promise and challenges}

The past decade has seen a surge of interest in replacing or augmenting classical closures with machine-learned models trained on high-fidelity data.
\citet{Ling2016} introduced the tensor basis neural network (TBNN), which embeds Galilean invariance into the network architecture by predicting coefficients of the Pope integrity basis \citep{Pope1975}.
\citet{Kaandorp2020} proposed the tensor basis random forest (TBRF), which achieves higher offline accuracy by using ensemble methods with the same invariant features and tensor basis.
\citet{Wu2018} explored physics-informed approaches that enforce realizability constraints on the predicted Reynolds stress tensor.

These data-driven closures have demonstrated improved \emph{a priori} accuracy on held-out datasets, particularly for flows with strong anisotropy, separation, and secondary motions where linear eddy viscosity models fail.
However, several challenges remain:
\begin{itemize}
    \item \textbf{A priori vs.\ a posteriori performance}: offline accuracy on held-out data does not guarantee improved predictions when the model is coupled into a CFD solver. The feedback loop between the closure and the mean-flow equations can amplify small errors.
    \item \textbf{Computational cost}: neural network inference involves matrix multiplies and nonlinear activation functions that differ fundamentally from the stencil-based operations of classical closures. The cost implications on modern GPU hardware are rarely quantified.
    \item \textbf{Generalization}: models trained on specific flow configurations may not transfer to geometries or Reynolds numbers outside the training distribution.
\end{itemize}

\subsection{The gap: cost is never measured}

Most existing evaluations of data-driven turbulence closures focus exclusively on prediction accuracy, either \emph{a priori} (offline, on held-out data) or, less commonly, \emph{a posteriori} (coupled into a solver).
A smaller body of work performs \emph{a posteriori} evaluation but rarely compares multiple architectures head-to-head in the same solver with the same training data.
Crucially, the \emph{computational cost} of neural network closures relative to classical models is almost never reported, and a detailed breakdown of \emph{where} the cost comes from---feature computation, matrix multiplies, activation functions, memory traffic---is entirely absent from the literature.

This omission matters because the cost--accuracy tradeoff is the fundamental decision a practitioner faces when selecting a turbulence model.
A model that is 10\% more accurate but 10$\times$ more expensive may be impractical for design optimization studies requiring thousands of solver evaluations.
Conversely, a model that adds only 5\% overhead while providing meaningfully better predictions for separated flows could be immediately useful.

\subsection{Contributions}

This work provides a systematic, apples-to-apples comparison of data-driven and classical turbulence closures, spanning the full spectrum from zero-equation algebraic models to random forests. Our specific contributions are:

\begin{enumerate}
    \item \textbf{Unified evaluation framework}: Five neural network / machine learning architectures (MLP, MLP-Large, TBNN, PI-TBNN, TBRF) and four classical models (baseline, $k$-$\omega$, SST, EARSM) are trained on the same dataset \citep{McConkey2021} with the same case-holdout protocol and implemented in a single GPU-accelerated incompressible Navier--Stokes solver.

    \item \textbf{Detailed computational anatomy}: We provide a per-operation breakdown of computational cost for each closure type, identifying the specific operations (matrix multiplies, \texttt{tanh} activations, tree traversals, transport equation solves) that dominate wall-clock time. This goes beyond aggregate timing to explain \emph{why} each model costs what it does.

    \item \textbf{Cost--accuracy Pareto analysis}: We map the tradeoff between prediction accuracy (\emph{a priori} and \emph{a posteriori}) and computational cost on both GPU and CPU hardware, identifying which models lie on the Pareto frontier and which are dominated.

    \item \textbf{Negative result on physics-informed penalties}: We show that realizability constraints imposed via loss-function penalties do not improve TBNN accuracy, because the tensor basis architecture already constrains outputs to be near-realizable.

    \item \textbf{TBRF solver integration}: We implement random forest inference in a CFD solver for the first time, demonstrating that the TBRF's superior offline accuracy comes at prohibitive online cost due to branch-heavy tree traversals and poor GPU cache utilization.
\end{enumerate}

The remainder of this paper is organized as follows.
Section~\ref{sec:solver} describes the solver architecture and GPU implementation.
Section~\ref{sec:closures} details the classical and data-driven closure models, including their computational cost per cell.
Section~\ref{sec:training} describes the training methodology and solver integration.
Section~\ref{sec:apriori} presents \emph{a priori} evaluation results.
Section~\ref{sec:aposteriori} presents \emph{a posteriori} results on canonical flow cases.
Section~\ref{sec:cost} provides the computational cost analysis and Pareto tradeoff.
Section~\ref{sec:discussion} discusses practical implications and limitations.
Section~\ref{sec:conclusions} summarizes our findings and recommendations.

\section{Solver Architecture}
\label{sec:solver}

\subsection{Governing equations}

We solve the incompressible Navier--Stokes equations,
\begin{align}
    \nabla \cdot \bm{u} &= 0, \label{eq:continuity} \\
    \frac{\partial \bm{u}}{\partial t} + (\bm{u} \cdot \nabla)\bm{u} &= -\nabla p + \nu \nabla^2 \bm{u} + \bm{f}, \label{eq:momentum}
\end{align}
where $\bm{u}$ is the velocity field, $p$ the kinematic pressure, $\nu$ the kinematic viscosity, and $\bm{f}$ a body force (e.g., a mean pressure gradient driving channel flow).
When a turbulence closure is active, the molecular viscosity $\nu$ is replaced by an effective viscosity $\nueff = \nu + \nut$, where $\nut$ is the eddy viscosity predicted by the closure model.
For anisotropic closures (EARSM, TBNN, TBRF), the Reynolds stress divergence is added directly as a source term rather than through an eddy viscosity.

\subsection{Fractional-step projection method}

Time integration follows the fractional-step (pressure-projection) method:
\begin{enumerate}
    \item \textbf{Turbulence update}: Advance the turbulence model to compute updated $\nut$ (or anisotropic stress tensor $\tau_{ij}$).
    \item \textbf{Predictor}: Compute an intermediate velocity $\bm{u}^*$ by advancing the momentum equation without the pressure gradient:
    \begin{equation}
        \bm{u}^* = \bm{u}^n + \Delta t \left( -\bm{C}(\bm{u}^n) + \bm{D}(\bm{u}^n) + \bm{f} \right),
    \end{equation}
    where $\bm{C}$ and $\bm{D}$ denote the convective and diffusive operators, respectively.
    Convection uses the skew-symmetric form for improved energy conservation.
    \item \textbf{Pressure Poisson}: Solve
    \begin{equation}
        \nabla^2 p' = \frac{1}{\Delta t} \nabla \cdot \bm{u}^*
    \end{equation}
    for the pressure correction $p'$.
    \item \textbf{Velocity correction}: Project to a divergence-free field:
    \begin{equation}
        \bm{u}^{n+1} = \bm{u}^* - \Delta t \, \nabla p'.
    \end{equation}
\end{enumerate}
This sequence is embedded within a strong stability preserving Runge--Kutta (SSP-RK3) time integrator, with steps 2--4 repeated at each stage with appropriate SSP weights.
Adaptive time stepping enforces directional CFL constraints: a relaxed $\mathrm{CFL}_{xz}$ for the streamwise and spanwise directions (where grid spacing is uniform) and a strict $\mathrm{CFL}_{\max}$ for the wall-normal direction (where the stretched grid produces small $\Delta y$ near walls).

\subsection{Staggered MAC grid}

Spatial discretization uses a staggered marker-and-cell (MAC) grid with second-order central differences.
Velocity components are stored at cell faces ($u$ at $x$-faces, $v$ at $y$-faces, $w$ at $z$-faces) and pressure at cell centers.
This arrangement guarantees pressure--velocity coupling without checkerboard oscillations and ensures that the discrete divergence and gradient operators compose to the exact discrete Laplacian:
\begin{equation}
    D \cdot G = L,
\end{equation}
a property critical for the projection method to enforce the incompressibility constraint~\eqref{eq:continuity} to machine precision.
For stretched grids in the wall-normal direction, this compatibility is maintained through precomputed metric arrays ($\Delta y_v$, $\Delta y_c$, $y_{\mathrm{Lap},*}$) that encode the non-uniform spacing in the discrete operators.

\subsection{GPU acceleration}

The solver achieves GPU acceleration through OpenMP target offload directives, maintaining a single source code path for both CPU and GPU execution.
All field data (velocity, pressure, turbulence quantities) is persistently mapped to GPU memory at initialization and remains on device throughout the simulation.
GPU--CPU data transfer occurs only during I/O operations (writing output files) and at initialization/finalization.

Key design decisions for GPU performance include:
\begin{itemize}
    \item \textbf{Single code path}: Arithmetic kernels contain no \texttt{\#ifdef} guards---the OpenMP target pragmas are silently ignored in CPU builds, ensuring identical numerical behavior on both platforms.
    \item \textbf{Free-function kernels}: GPU compute kernels are implemented as free functions rather than class methods to avoid implicit transfer of the \texttt{this} pointer (an nvc++ compiler requirement).
    \item \textbf{Persistent device memory}: The \texttt{OmpDeviceBuffer} wrapper manages GPU allocations with RAII semantics. Data is mapped once and remains on device.
    \item \textbf{Reduction-based diagnostics}: Scalar diagnostics (kinetic energy, maximum velocity, divergence $L_\infty$ norm) are computed via GPU reductions without transferring full fields to the host.
\end{itemize}

The solver is compiled with the NVHPC compiler (nvc++) for GPU builds and GCC for CPU-only builds.

\subsection{Poisson solver}

The pressure Poisson equation is solved using a fast Fourier transform (FFT) approach via the cuFFT library.
For directions with periodic boundary conditions, standard FFTs are applied; for directions with Dirichlet or Neumann conditions, discrete cosine/sine transforms are used.
The FFT approach reduces the 3D Poisson problem to a set of independent tridiagonal systems in the wall-normal direction, which are solved directly.
For stretched grids, the tridiagonal coefficients incorporate the mesh metric arrays to maintain $D \cdot G = L$ consistency, and the right-hand side uses volume-weighted mean subtraction to satisfy solvability.

The FFT Poisson solver achieves 22--32$\times$ speedup over the geometric multigrid alternative on GPU for stretched grids, making it the preferred backend for all production runs.

\subsection{Immersed boundary method}

Complex geometries (cylinder, airfoil, periodic hills) are handled via an immersed boundary method (IBM) using direct forcing.
At initialization, signed distance functions define the solid boundary, and weight arrays are precomputed: $w = 0$ inside the solid, $w = 1$ in the fluid, with a smooth transition in the forcing region.
During each time step, the IBM applies element-wise multiplication of the velocity field by these weight arrays---a trivially parallel operation that adds less than 0.3\% overhead on GPU.
The Poisson equation right-hand side is masked at solid cells via a separate GPU kernel, avoiding any CPU--GPU synchronization.

\subsubsection{Ghost-cell immersed boundary method}
\label{sec:ghost_cell_ibm}

The primary IBM approach uses the ghost-cell method of \citet{Fadlun2000}.
For each ghost cell (the first solid cell adjacent to the fluid--solid interface), the velocity is reconstructed by interpolating from the fluid side along the surface-normal direction to enforce the no-slip condition at the boundary.
This yields sharp boundary representation with second-order accuracy in the boundary location, avoiding the $O(\eta)$ accuracy limitation of volume penalization methods~\citep{Angot1999}.

The solver also retains a semi-implicit volume penalization option,
\begin{equation}
    \bm{u} \leftarrow \bm{u} \left[ 1 - (1 - w) \min\!\left(\frac{\Delta t}{\eta}, 1\right) \right],
    \label{eq:volume_penalization}
\end{equation}
where $\eta$ is a penalization timescale and $w$ is the solid mask ($w=0$ inside the body, $w=1$ in the fluid).
This penalization is used during the SST warm-up phase for the periodic hills geometry, where it provides robust stabilization of the initial transient; the solver then switches to the ghost-cell method for the production measurement phase.
For cylinder and sphere cases, ghost-cell forcing is used throughout.

In all cases, the re-forcing step after the pressure correction is removed: only the pre-correction forcing (step~5 in Section~\ref{sec:solver}) and the Poisson RHS masking (step~7) are retained.
This eliminates a feedback instability that arises when the corrected velocity is re-forced, generating divergence that is never projected out.

\subsubsection{Grid convergence of the ghost-cell IBM}
\label{sec:ibm_grid_convergence}

The ghost-cell method of \citet{Fadlun2000} is first-order accurate in the grid spacing $h$, since the velocity at each ghost cell is reconstructed via linear interpolation along the surface normal.
We verify this convergence behavior on two canonical cases: the circular cylinder at $\Rey = 100$ and the sphere at $\Rey = 200$.

Table~\ref{tab:ibm_convergence_cyl} shows the cylinder drag coefficient $C_D$ as a function of grid resolution, expressed as cells per diameter $D/h$.
At 8 cells/$D$ the error exceeds 50\%, but by 24 cells/$D$ the ghost-cell method recovers $C_D$ to within 9\% of the reference value $C_D = 1.35$.
Adding SST warm-up (Section~\ref{sec:sst_warmup}) further improves the result to $C_D = 1.40$ (3.6\% error), indicating that the turbulence initialization contributes meaningfully once the IBM resolution is adequate.

\begin{table}[htbp]
\centering
\caption{Grid convergence of ghost-cell IBM for circular cylinder at $\Rey = 100$. Reference $C_D = 1.35$.}
\label{tab:ibm_convergence_cyl}
\begin{tabular}{lrrr}
\toprule
Grid & Cells/$D$ & $C_D$ & Error \\
\midrule
$128 \times 96$  & 8  & 0.63 & 53\% \\
$384 \times 288$ & 24 & 1.47 & 9\%  \\
$384 \times 288$ (SST warm-up) & 24 & 1.40 & 3.6\% \\
\bottomrule
\end{tabular}
\end{table}

Table~\ref{tab:ibm_convergence_sph} presents the corresponding data for the sphere.
Convergence is substantially slower: at 24 cells/$D$, the error remains 51\%.
This is a known property of first-order immersed boundary methods on staggered grids---the three-dimensional interpolation geometry at curved surfaces is more complex than in two dimensions, and the volume fraction of cells cut by the interface scales as $O(h^{-2})$ in 3D versus $O(h^{-1})$ in 2D.
Notably, the baseline simulation (no SST warm-up) produces the same $C_D$ as the SST-initialized run at $384 \times 256^2$, confirming that the IBM grid resolution is the convergence bottleneck rather than the turbulence model.

\begin{table}[htbp]
\centering
\caption{Grid convergence of ghost-cell IBM for sphere at $\Rey = 200$. Reference $C_D = 0.77$.}
\label{tab:ibm_convergence_sph}
\begin{tabular}{lrrr}
\toprule
Grid & Cells/$D$ & $C_D$ & Error \\
\midrule
$128 \times 96^2$  & 10 & 0.17 & 78\% \\
$192 \times 128^2$ & 13 & 0.25 & 68\% \\
$256 \times 192^2$ & 17 & 0.21 & 73\% \\
$384 \times 256^2$ & 24 & 0.38 & 51\% \\
\bottomrule
\end{tabular}
\end{table}

These results are consistent with the $O(h)$ convergence rate reported in the IBM literature \citep{Fadlun2000}.
For the cylinder, 24 cells/$D$ provides sufficient resolution for quantitative $C_D$ comparison (${<}10\%$ error).
For the sphere, while absolute $C_D$ values have not converged at affordable resolutions, relative comparisons between turbulence models at a fixed grid resolution remain meaningful---any model-to-model differences are attributable to the closure rather than to IBM discretization error.

\subsubsection{Grid resolution requirements for IBM stability}
\label{sec:ibm_grid_requirements}

The combination of the explicit RK3 time integrator with the ghost-cell IBM imposes minimum grid resolution requirements that are substantially higher than those for body-fitted solvers.
The IBM introduces a pressure transient at the fluid--solid interface during the initial time steps: the penalization term damps the body-interior velocity toward zero, creating a localized divergence source that the Poisson solver must correct.
With insufficient resolution, this transient produces oscillations that grow exponentially and diverge within a few hundred steps.

Table~\ref{tab:ibm_min_grid} reports the minimum stable grid for each a~posteriori case, determined empirically by testing progressively finer grids until long-time stability is achieved (${>}2000$ steps without divergence).
The periodic hills geometry requires approximately $4\times$ the resolution that would be sufficient for a body-fitted solver, while the sphere requires $2\times$ refinement in each direction.
This overhead is a direct consequence of the explicit time integrator: implicit solvers (e.g., OpenFOAM's SIMPLE algorithm, commonly used for RANS periodic hills at Re${=}5{,}600$) can accommodate the IBM pressure transient without excessive resolution.

\begin{table}[htbp]
\centering
\caption{Minimum stable grid resolution for each IBM case with the explicit RK3 solver. The duct uses wall-fitted boundaries (no IBM). The ``factor'' column indicates the grid refinement relative to an initial estimate based on cell-count matching with body-fitted RANS literature.}
\label{tab:ibm_min_grid}
\begin{tabular}{lrrrr}
\toprule
Case & Grid & Cells & Cells/$D$ & Factor \\
\midrule
Cylinder Re${=}100$ & $384 \times 288$ & 0.1M & 24 & $1\times$ \\
Hills Re${=}5{,}600$ & $1{,}536 \times 768$ & 1.2M & --- & $4\times$ \\
Duct Re$_b{=}3{,}500$ & $96^3$ & 0.9M & --- & $1\times$ \\
Sphere Re${=}200$ & $384 \times 256^2$ & 25.2M & 48 & $2\times$/dir \\
\bottomrule
\end{tabular}
\end{table}

The volume penalization (Eq.~\ref{eq:volume_penalization}) with $\eta = 0.1$ is applied during the SST warm-up phase for all IBM cases to stabilize the initial transient.
After warm-up, the solver switches to ghost-cell forcing for higher accuracy.

\subsection{SST warm-up initialization}
\label{sec:sst_warmup}

Non-transport turbulence models---including mixing-length, GEP, and all neural network closures---compute the eddy viscosity $\nut$ from the local velocity gradient $\partial u_i / \partial x_j$.
When the simulation is initialized from a uniform velocity field, all velocity gradients are zero and these models produce $\nut = 0$, causing the early-time flow development to proceed as if laminar.
This can delay or entirely prevent the establishment of a turbulent mean flow.

To address this, the solver implements a two-phase initialization:
\begin{enumerate}
    \item \textbf{SST warm-up}: The Menter SST $k$--$\omega$ transport model~\citep{Menter1994} is run for a specified physical time (\texttt{warmup\_time}), developing the velocity, $k$, and $\omega$ fields from the uniform initial condition. Because the SST model integrates transport equations for $k$ and $\omega$, it produces nonzero $\nut$ even from uniform flow via the production terms.
    \item \textbf{Model switch}: The target turbulence closure replaces SST and operates on the developed flow field. The $k$ and $\omega$ fields are discarded (unless the target model is itself a transport model).
\end{enumerate}
Timer statistics are reset at the model switch so that reported wall-clock times reflect only the cost of the target closure, not the SST warm-up phase.

\subsection{Three-phase timing methodology}
\label{sec:timing_methodology}

Rigorous cost measurement requires separating model startup transients from steady-state operation.
Each model evaluation uses a three-phase protocol:
\begin{enumerate}
    \item \textbf{SST warm-up} (Section~\ref{sec:sst_warmup}): Develop the flow field using SST $k$--$\omega$ for a specified physical time. This phase is not timed.
    \item \textbf{Model stabilization}: The target closure runs for a specified number of warm-up steps (\texttt{warmup\_steps}) to allow its $\nut$ field to equilibrate with the developed flow. This phase is not timed.
    \item \textbf{Measurement}: Timing is recorded over the remaining steps. Only this phase contributes to the reported wall-clock cost.
\end{enumerate}
All three phases use adaptive time stepping with identical CFL constraints, so the time step adjusts naturally to the effective viscosity of each model.
For the fixed-$\Delta t$ cost comparisons (Section~\ref{sec:cost}), a common time step is imposed across all models to ensure that the Poisson, convection, and diffusion costs are identical and only the turbulence closure cost varies.
The relevant configuration parameters are \texttt{warmup\_model}, \texttt{warmup\_time}, and \texttt{warmup\_steps}.

\subsection{Performance validation}
\label{sec:perf_validation}

To establish that our solver is well-optimized---and therefore that the turbulence closure cost measurements reported in Section~\ref{sec:cost} are meaningful---we benchmark against CaNS \citep{Costa2018} and CaLES \citep{Xiao2024}, widely used GPU-accelerated incompressible Navier--Stokes solvers.
CaNS employs the same numerical approach as our solver: second-order staggered finite differences, SSP-RK3 time integration, and an FFT-based direct Poisson solver.
It is written in Fortran with CUDA Fortran (CUF) kernels and uses the cuDecomp library for multi-GPU pencil decomposition.
CaLES extends CaNS with LES subgrid-scale models and reports single-GPU performance on A100.

We match the CaNS benchmark configuration exactly: channel flow at $\Rey_\tau \approx 590$ on a $512 \times 256 \times 144$ grid (18.9M cells), with second-order central differences, explicit RK3 time integration, explicit diffusion, and FFT-based Poisson solver.

\begin{table}[htbp]
\centering
\caption{Solver throughput comparison with CaNS \citep{Costa2018} and CaLES \citep{Xiao2024} on identical discretization (2nd-order central FD, RK3, explicit diffusion, FFT Poisson). CaNS grid: $512 \times 256 \times 144$ (18.9M cells), V100-SXM2-32GB on DGX-2. CaLES grid: $512 \times 384 \times 1440$ (283M cells), single A100-SXM4-64GB.}
\label{tab:solver_throughput}
\begin{tabular}{llrrr}
\toprule
Solver & GPU configuration & ms/step & Mcells/s & Mcells/s/GPU \\
\midrule
This work   & $1 \times$ V100-PCIe-16GB  & 74.5  & 253  & 253 \\
This work   & $1 \times$ H100-SXM5-80GB  & ---   & ---  & --- \\
\midrule
CaNS        & $4 \times$ V100-SXM2-32GB  & 481    & 39   & 9.8 \\
CaNS        & $16 \times$ V100-SXM2-32GB & 140    & 135  & 8.4 \\
\bottomrule
\end{tabular}
\end{table}

Table~\ref{tab:solver_throughput} shows that our solver achieves 253~Mcells/s on a single V100 at 16.8M cells, which is $6.5\times$ the total throughput of CaNS running on 4~V100s ($39$~Mcells/s) at a comparable grid size, and $1.9\times$ CaNS on 16~V100s ($135$~Mcells/s).
The per-GPU efficiency is $26\times$ higher (253 vs.\ 9.8~Mcells/s/GPU).

The large per-GPU advantage over CaNS's multi-GPU configuration arises from the distributed FFT communication overhead inherent in CaNS's pencil decomposition.
CaNS uses cuDecomp for 2D pencil transposes that require global all-to-all communication.
With RK3, three Poisson solves are performed per time step, each requiring forward and backward FFTs with multiple transposes.
\citet{Costa2018} report that the Poisson solver consumes 60--70\% of wall-clock time in the multi-GPU configuration, with much of that attributable to transposes.
Our single-GPU implementation avoids this overhead entirely, performing the full FFT on-chip via cuFFT.

\subsubsection{Profiling-guided optimization}

To close the remaining gap with CaLES \citep{Xiao2024}---a single-GPU LES code built on CaNS---we used NVIDIA Nsight Systems (nsys) to profile our solver at the kernel level.
The initial profile on V100 at $256^3$ revealed that the Poisson solver consumed 73\% of kernel time, compared to CaLES's 38\%.

The root cause was identified by examining the cuFFT plan configuration.
Our batched 2D FFT used a strided memory layout (\texttt{istride}=$N_y$, \texttt{idist}=1), meaning consecutive elements in the transform dimensions ($x$, $z$) were separated by $N_y = 256$ doubles (2~KB) in memory.
This non-unit stride caused non-coalesced memory access within cuFFT's internal radix kernels: each warp of 32~threads accessed 32 different cache lines instead of 1--2 contiguous lines.
A standalone cuFFT benchmark confirmed an $11\times$ penalty: 8.25~ms/call with stride=$N_y$ versus 0.74~ms/call with stride=1 for identical transforms.

The fix required reorganizing the packed data layout from $\text{packed}[(i N_z + k) N_y + j]$ (y-interleaved, stride=$N_y$ for FFT) to $\text{packed}[j (N_x N_z) + i N_z + k]$ (y-batched, stride=1 for FFT).
Since the cuSPARSE batched tridiagonal solver requires mode-major layout ($\text{hat}[m N_y + j]$), we added a shared-memory tiled matrix transpose (32$\times$32 tiles with bank-conflict-free padding) before and after the tridiagonal solve.
The transpose cost ($\sim$0.3~ms/call) is negligible compared to the cuFFT savings.

A second optimization separated the post-FFT data unpacking and boundary condition application from a single kernel with scattered writes into two kernels: a pure transpose (output-coalesced, $4.7\times$ faster) and a separate ghost cell fill (no branch divergence).

Table~\ref{tab:optimization} summarizes the cumulative impact.

\begin{table}[htbp]
\centering
\caption{Profiling-guided optimizations and their impact on total step time (V100, $256^3$, channel flow, RK3).}
\label{tab:optimization}
\begin{tabular}{lrrr}
\toprule
Version & ms/step & ns/point & Cumulative speedup \\
\midrule
Initial implementation                 & 141.0 & 8.40 & $1.00\times$ \\
+ Output-coalesced unpack kernel       & 123.9 & 7.39 & $1.14\times$ \\
+ Contiguous cuFFT layout (stride=1)   & 74.5  & 4.44 & $1.89\times$ \\
\bottomrule
\end{tabular}
\end{table}

After optimization, the kernel-level breakdown (nsys) shows a well-balanced solver with no single dominant bottleneck (Table~\ref{tab:kernel_breakdown}).

\begin{table}[htbp]
\centering
\caption{GPU kernel time breakdown after optimization (V100, $256^3$, 20 steps). Percentages are of total kernel time.}
\label{tab:kernel_breakdown}
\begin{tabular}{lrr}
\toprule
Category & ms/step & \% of kernels \\
\midrule
Poisson solver (FFT + tridiag + pack/unpack) & 35.0 & 30\% \\
\quad cuFFT transforms                       & 16.6 & 14\% \\
\quad cuSPARSE tridiagonal                   & 10.5 & 9\% \\
\quad Pack/unpack + transpose                & 7.9  & 7\% \\
RK3 stepping (blend + substep)               & 33.5 & 28\% \\
Convection + diffusion                       & 17.3 & 15\% \\
Velocity correction + projection             & 18.3 & 16\% \\
Divergence                                   & 6.8  & 6\% \\
Boundary conditions                          & 2.3  & 2\% \\
Other                                        & 3.5  & 3\% \\
\bottomrule
\end{tabular}
\end{table}

We also evaluated replacing the cuSPARSE PCR-based batched tridiagonal solver with a custom Thomas algorithm (as used by CaNS), but found cuSPARSE to be faster at $N_y = 256$: the PCR algorithm's $O(\log N)$ parallel steps outperform Thomas's $O(N)$ serial sweep when sufficient inter-system parallelism exists, and the Thomas solver's per-thread local array ($N_y$ doubles) causes register spilling to local memory on GPU.

\subsubsection{Comparison with CaLES}

CaLES \citep{Xiao2024} reports a per-subroutine timing breakdown for its base solver (DNS mode, no turbulence model) on a single A100-SXM4-64GB GPU with a $512 \times 384 \times 1440$ grid (283M cells):
Poisson solver 38\%, RHS (convection + diffusion) 30\%, velocity correction 13\%, and SGS model 19\% (Smagorinsky, not applicable to our RANS comparison).
Their DNS throughput is 1.50~ns/step/point ($\sim$667~Mcells/s).

Our optimized solver achieves 4.44~ns/step/point on V100 at $256^3$.
Accounting for the V100/A100 memory bandwidth ratio (0.90 vs.\ 2.0~TB/s), the bandwidth-normalized per-point cost is within a factor of 1.3 of CaLES.
The residual gap is attributable to (i) smaller grid size (16.8M vs.\ 283M cells, reducing GPU occupancy), (ii) programming model differences (OpenMP target offload vs.\ CUDA Fortran/OpenACC), and (iii) CaNS's in-place cuFFT execution which avoids one buffer allocation.

These results confirm that our solver is competitive with established, heavily optimized GPU codes, and that the turbulence closure cost measurements reported in Section~\ref{sec:cost} are measured against a well-optimized baseline---not inflated by solver inefficiencies.

\section{Turbulence Closure Models}
\label{sec:closures}

We evaluate nine turbulence closure models spanning the full hierarchy from zero-equation algebraic to random forest.
All models are implemented in the same solver and share the same spatial discretization, time integration, and Poisson solver.
Table~\ref{tab:model_summary} summarizes their key properties.

\begin{table}[htbp]
\centering
\caption{Summary of turbulence closure models evaluated. FLOPs/cell counts are for the turbulence update only (excluding shared solver operations). Weight sizes refer to the model parameters stored on device.}
\label{tab:model_summary}
\begin{tabular}{lllrrl}
\toprule
Model & Type & Architecture & Parameters & FLOPs/cell & Weights \\
\midrule
Baseline       & Algebraic    & Mixing length              & 0       & $\sim$2    & --- \\
$k$-$\omega$   & Transport    & 2-equation                 & 0       & $\sim$50   & --- \\
SST            & Transport    & 2-equation + blending      & 0       & $\sim$80   & --- \\
EARSM          & Algebraic RS & SST + tensor algebra       & 0       & $\sim$120  & --- \\
GEP            & Algebraic RS & Symbolic regression        & 0       & $\sim$100  & --- \\
MLP            & NN (scalar)  & $5 \to 32 \to 32 \to 1$    & 1{,}249  & $\sim$1{,}249 & \SI{56}{KB} \\
MLP-Large      & NN (scalar)  & $5 \to 128^4 \to 1$        & 50{,}049 & $\sim$50{,}049 & \SI{896}{KB} \\
TBNN           & NN (tensor)  & $5 \to 64^3 \to 10$        & 9{,}354  & $\sim$10{,}300 & \SI{196}{KB} \\
PI-TBNN        & NN (tensor)  & Same as TBNN               & 9{,}354  & $\sim$10{,}300 & \SI{196}{KB} \\
TBRF           & RF (tensor)  & 10 coeff.\ $\times$ 10 trees & $\sim$2.8M nodes & $\sim$2{,}000 branches & \SI{56}{MB} \\
\bottomrule
\end{tabular}
\end{table}

\subsection{Classical models}

\subsubsection{Baseline mixing length}

The algebraic baseline computes the eddy viscosity from a mixing length $\ell_{\mathrm{mix}}$ and the local velocity gradient:
\begin{equation}
    \nut = \ell_{\mathrm{mix}}^2 \left| \frac{\partial u}{\partial y} \right|.
\end{equation}
This requires one division and one multiply per cell, reading 3 velocity fields and writing 1 eddy viscosity field.
On GPU, this is a memory-bandwidth-bound operation with perfect parallelism.

\subsubsection{$k$-$\omega$ transport model}

The $k$-$\omega$ model \citep{Wilcox1988} solves transport equations for the turbulent kinetic energy $k$ and the specific dissipation rate $\omega$:
\begin{align}
    \frac{\partial k}{\partial t} + u_j \frac{\partial k}{\partial x_j} &= P_k - \beta^* k \omega + \frac{\partial}{\partial x_j}\left[(\nu + \sigma_k \nut)\frac{\partial k}{\partial x_j}\right], \\
    \frac{\partial \omega}{\partial t} + u_j \frac{\partial \omega}{\partial x_j} &= \alpha \frac{\omega}{k} P_k - \beta \omega^2 + \frac{\partial}{\partial x_j}\left[(\nu + \sigma_\omega \nut)\frac{\partial \omega}{\partial x_j}\right].
\end{align}
Each equation involves convective and diffusive flux computation (stencil operations) plus source and sink terms, totaling approximately 50~FLOPs per cell.
Destruction terms ($\beta^* k \omega$, $\beta \omega^2$) are treated point-implicitly to avoid stiffness-induced blow-up at wall cells, where $\omega_{\mathrm{wall}} \sim O(10^5)$.

\subsubsection{SST $k$-$\omega$}

The SST model \citep{Menter1994} augments $k$-$\omega$ with blending functions $F_1$ and $F_2$ that transition between near-wall ($k$-$\omega$) and free-stream ($k$-$\varepsilon$) formulations, plus a cross-diffusion term $\mathrm{CD}_{k\omega}$.
Each blending function requires wall distance, $\nabla k \cdot \nabla \omega$, and several auxiliary quantities, raising the cost to approximately 80~FLOPs per cell (2.5$\times$ the base $k$-$\omega$ model).

\subsubsection{EARSM}

The explicit algebraic Reynolds stress model extends SST by computing the anisotropy tensor $a_{ij}$ from an algebraic relation involving the normalized strain rate $\Shat_{ij}$ and rotation rate $\Omhat_{ij}$ tensors:
\begin{equation}
    a_{ij} = f(\Shat, \Omhat),
\end{equation}
where the functional form follows Pope, Wallin--Johansson, or Gatski--Speziale closures.
The tensor algebra (3$\times$3 matrix multiplications, traces, invariant computation) adds approximately 40~FLOPs per cell on top of the SST transport, totaling $\sim$120~FLOPs per cell.
Despite the higher per-cell cost, EARSM retains the stencil-based memory access pattern of transport models and remains memory-bandwidth-bound on GPU.

\subsubsection{GEP}

The gene expression programming model \citep{Weatheritt2016} uses a symbolic expression discovered via genetic programming to compute the Reynolds stress anisotropy.
The expression is algebraic (no transport equations beyond SST), with a cost of $\sim$100~FLOPs per cell.

\subsection{Neural network architectures}

All neural network models take the same 5 scalar invariants of the non-dimensionalized strain rate and rotation rate tensors as input (Section~\ref{sec:training}).
The architectures differ in capacity, output dimensionality, and how the output maps to a turbulence closure.

\subsubsection{MLP (scalar closure)}

The multilayer perceptron predicts a scalar proxy for the eddy viscosity magnitude:
\begin{equation}
    |b| = f_{\mathrm{MLP}}(\hat{\lambda}_1, \ldots, \hat{\lambda}_5),
\end{equation}
with architecture $5 \to 32 \to 32 \to 1$ (1{,}249 parameters).
Hidden layers use $\tanh$ activations; the output layer is linear.
The inference cost per cell is:
\begin{itemize}
    \item Layer 0 ($5 \to 32$): 160 multiplies + 32 adds + 32 $\tanh$
    \item Layer 1 ($32 \to 32$): 1{,}024 multiplies + 32 adds + 32 $\tanh$
    \item Layer 2 ($32 \to 1$): 32 multiplies + 1 add
    \item \textbf{Total}: 1{,}249 FLOPs + 64 $\tanh$ evaluations
\end{itemize}
The weight matrix (\SI{56}{KB}) fits in GPU L1 cache, enabling broadcast to all threads with minimal memory traffic.

\subsubsection{MLP-Large (scalar closure)}

A larger MLP with architecture $5 \to 128 \to 128 \to 128 \to 128 \to 1$ (50{,}049 parameters) serves as a capacity scaling test.
The dominant cost is four $128 \times 128$ matrix-vector products (16{,}384 multiplies each), totaling $\sim$50{,}000~FLOPs per cell with 512~$\tanh$ evaluations.
The weight matrix (\SI{896}{KB}) exceeds L1 but fits in L2 cache.

\subsubsection{TBNN (tensor basis closure)}
\label{sec:tbnn}

The tensor basis neural network \citep{Ling2016} predicts 10 scalar coefficients $g^{(n)}$ that are contracted with the Pope integrity basis \citep{Pope1975} to reconstruct the anisotropy tensor:
\begin{equation}
    b_{ij} = \sum_{n=1}^{10} g^{(n)}(\hat{\lambda}_1, \ldots, \hat{\lambda}_5) \, T^{(n)}_{ij},
    \label{eq:tbnn}
\end{equation}
where $T^{(n)}$ are the 10 symmetric, traceless tensor basis functions of $\Shat$ and $\Omhat$.
The architecture is $5 \to 64 \to 64 \to 64 \to 10$ (9{,}354 parameters) with $\tanh$ hidden activations.

The per-cell computational cost breaks down as follows:
\begin{itemize}
    \item \textbf{Feature computation}: velocity gradients (stencil), $S_{ij}$ and $\Omega_{ij}$ (symmetrize/antisymmetrize), non-dimensionalize by $k/\varepsilon$, compute 5 invariants (matrix products + traces)---shared with EARSM, $\sim$80 FLOPs.
    \item \textbf{Tensor basis computation}: 10 basis tensors $T^{(n)}$ from products of $\Shat$ and $\Omhat$, each requiring one or more $3 \times 3$ matrix multiplies---$\sim$900~FLOPs.
    \item \textbf{NN forward pass}: 4 layers totaling 9{,}354 FLOPs + 192 $\tanh$ evaluations. The cost is dominated by the two $64 \times 64$ hidden layers (4{,}096 FLOPs each, comprising 87\% of the network's compute).
    \item \textbf{Tensor contraction}: $b_{ij} = \sum_n g^{(n)} T^{(n)}_{ij}$---60 multiplies + 60 adds.
    \item \textbf{Total}: $\sim$10{,}300 FLOPs/cell, approximately 7.5$\times$ the MLP.
\end{itemize}

The TBNN weight matrix (\SI{196}{KB}) fits in GPU L2 cache.
The architecture guarantees Galilean invariance, symmetry, and zero trace of the predicted $b_{ij}$ by construction.

\subsubsection{PI-TBNN (physics-informed TBNN)}

The PI-TBNN shares the same architecture as the TBNN but is trained with an augmented loss function that penalizes violations of the realizability constraints $b_{ii} \geq -1/3$ and $b_{ii} \leq 2/3$ (see Section~\ref{sec:training}).
At inference time, the PI-TBNN is computationally identical to the TBNN.

\subsubsection{TBRF (tensor basis random forest)}

The tensor basis random forest \citep{Kaandorp2020} replaces the neural network with an ensemble of decision trees.
For each of the 10 basis coefficients $g^{(n)}$, a separate random forest of 10--200 trees (depth 20) is trained on least-squares-extracted coefficients.

The inference cost per cell differs fundamentally from neural networks:
\begin{itemize}
    \item For each cell, for each of 10 coefficients, traverse 10 trees to depth $\sim$20: this requires $10 \times 10 \times 20 = 2{,}000$ comparison-and-branch operations.
    \item Each branch decision accesses a different memory location in the tree structure (feature index, threshold, child pointers), producing scattered, non-coalesced memory accesses.
    \item The 10-tree TBRF requires \SI{56}{MB} of tree data, which does not fit in any GPU cache level (L1: \SI{128}{KB}, L2: \SI{48}{MB} on A100).
    \item Warp divergence is severe: cells in the same warp follow different paths through each tree, serializing execution.
\end{itemize}

Contrast this with neural network inference, where every cell performs the \emph{same} matrix-vector multiply with the \emph{same} weights---a perfectly regular, data-parallel operation with high arithmetic intensity and coalesced memory access.

\subsection{Computational anatomy: why NNs are expensive on GPU}
\label{sec:comp_anatomy}

The fundamental mismatch between neural network inference and GPU architecture manifests in three ways:

\paragraph{Layer-sequential execution.}
Each MLP layer depends on the previous layer's output.
For a TBNN with 4 layers, this means 4 serial phases within each kernel.
Parallelism is only over (cell, neuron) pairs; the depth dimension is sequential.
This limits the arithmetic intensity compared to stencil operations that process all cells independently.

\paragraph{Activation functions.}
The $\tanh$ function costs approximately 20 GPU cycles per evaluation versus $\sim$4 cycles for a multiply.
With 192 $\tanh$ evaluations per cell for the TBNN, activation functions account for a significant fraction of the total compute.

\paragraph{Memory access patterns.}
Classical stencil-based closures (transport models) access compact neighborhoods of large arrays with perfect spatial locality and coalesced GPU memory access.
Neural network weights must be broadcast to all threads---this is efficient for small models (MLP, \SI{56}{KB}, fits in L1) but increasingly costly as model size grows.
The TBRF's tree data (\SI{56}{MB}) causes cache thrashing and random memory accesses, making it fundamentally poorly suited for GPU execution.

\section{Training and Solver Integration}
\label{sec:training}

\subsection{Dataset}

All models are trained on the curated turbulence modelling dataset of \citet{McConkey2021}, which contains 902{,}601 spatial points across 38 flow cases in four geometries: square duct (14 Reynolds numbers), periodic hills (5 slope angles), converging-diverging channel, and curved backward-facing step.
Each point provides RANS fields from a $k$-$\omega$ SST \citep{Menter1994} simulation (velocity gradients, $k$, $\varepsilon$, $\omega$, $\nut$) and DNS/LES reference data (Reynolds stress anisotropy tensor $b_{ij}$).

An important characteristic of this dataset: the square duct is the only genuinely three-dimensional flow.
Its DNS reference data \citep{Pinelli2010} contains nonzero cross-plane Reynolds stresses ($\overline{u'w'}$, $\overline{v'w'}$, i.e., $b_{13}$ and $b_{23}$) that drive Prandtl's secondary flow of the second kind.
However, the RANS input fields for the duct are effectively two-dimensional: the linear eddy viscosity model ($k$-$\omega$ SST) cannot predict the secondary vortices, so the RANS cross-plane velocities $V$ and $W$ remain near zero \citep{McConkey2021}.
The remaining three geometries (hills, CDC, CBFS) have homogeneous spanwise directions and strictly two-dimensional mean flow.
The models therefore learn the 3D anisotropy correction ($b_{13}$, $b_{23}$) from the duct cases while receiving near-zero RANS cross-plane velocity gradient inputs---the correction is driven entirely by the Pope invariants and tensor basis, not by RANS-predicted secondary flow.

\subsection{Feature engineering}
\label{sec:features}

\subsubsection{Input features: 5 Pope invariants}

Following \citet{Ling2016} and \citet{Pope1975}, the inputs to all neural network models are 5 scalar invariants of the non-dimensionalized mean strain rate and rotation rate tensors.

The strain and rotation rate tensors are computed from the RANS velocity gradients:
\begin{equation}
    S_{ij} = \frac{1}{2}\left(\frac{\partial U_i}{\partial x_j} + \frac{\partial U_j}{\partial x_i}\right), \quad
    \Omega_{ij} = \frac{1}{2}\left(\frac{\partial U_i}{\partial x_j} - \frac{\partial U_j}{\partial x_i}\right),
\end{equation}
and non-dimensionalized using the turbulence time scale:
\begin{equation}
    \Shat_{ij} = \frac{k}{\varepsilon} S_{ij}, \quad \Omhat_{ij} = \frac{k}{\varepsilon} \Omega_{ij},
\end{equation}
where $k$ and $\varepsilon$ are clamped to a minimum of $10^{-30}$ to prevent division by zero.
The 5 scalar invariants are:
\begin{equation}
    \lambda_1 = \mathrm{tr}(\Shat^2), \quad
    \lambda_2 = \mathrm{tr}(\Omhat^2), \quad
    \lambda_3 = \mathrm{tr}(\Shat^3), \quad
    \lambda_4 = \mathrm{tr}(\Omhat^2 \Shat), \quad
    \lambda_5 = \mathrm{tr}(\Omhat^2 \Shat^2).
\end{equation}
These form a complete integrity basis for the independent invariants of a symmetric and an antisymmetric $3 \times 3$ tensor \citep{Pope1975}.

Input features are Z-score standardized using means and standard deviations computed from the training set only.

\subsubsection{Tensor basis}

For models predicting the full anisotropy tensor (TBNN, PI-TBNN, TBRF), the 10-tensor Pope integrity basis $T^{(n)}_{ij}$ is computed from $\Shat$ and $\Omhat$ (see Appendix~\ref{sec:appendix_basis}).
Each basis tensor is symmetric and traceless, stored as 6 independent components.

\subsubsection{Training labels}

\paragraph{TBNN / PI-TBNN / TBRF.} The full anisotropy tensor $b_{ij} = \overline{u'_i u'_j} / (2k) - \delta_{ij}/3$ from DNS, represented as 6 symmetric components.
The network predicts the raw DNS anisotropy---not the discrepancy $\Delta b = b_{\mathrm{DNS}} - b_{\mathrm{RANS}}$---following \citet{Ling2016}.

\paragraph{MLP / MLP-Large.} A scalar proxy: the Frobenius-like norm $|b| = \sqrt{\sum_c b_c^2}$ computed over the 6 stored components.

\subsection{Case-holdout split}

We adopt a case-holdout protocol where entire flow cases are held out, testing geometric generalization rather than interpolation within a single flow configuration:

\begin{itemize}
    \item \textbf{Training} (18 cases, 271{,}924 points): Square duct at 14 Reynolds numbers, periodic hills at $\alpha = \{0.5, 1.0, 1.5\}$, converging-diverging channel.
    \item \textbf{Validation} (2 cases, 23{,}967 points): Square duct $\Rey = 2000$ (interpolation), periodic hills $\alpha = 0.8$ (interpolation).
    \item \textbf{Test} (2 cases, 51{,}844 points): Periodic hills $\alpha = 1.2$ (interpolation at new angle), curved backward-facing step $\Rey = 13{,}700$ (entirely new geometry).
\end{itemize}

\subsection{Training hyperparameters}

All neural networks are trained with the Adam optimizer \citep{Kingma2015} using cosine annealing with warm restarts (initial period $T_0 = 200$ epochs, period multiplier $T_{\mathrm{mult}} = 2$), initial learning rate $10^{-3}$, minimum learning rate $10^{-6}$, and batch size 256.
Training runs for up to 1{,}000 epochs with early stopping (patience 150 epochs on validation MSE).
The best checkpoint (lowest validation MSE) is selected as the final model.
No regularization (L2 weight decay, dropout) is applied.
Full training details are provided in Appendix~\ref{sec:appendix_training}.

\subsection{PI-TBNN realizability penalty}

The PI-TBNN augments the standard MSE loss with a soft realizability penalty:
\begin{equation}
    \mathcal{L} = \mathrm{MSE}(b_{\mathrm{pred}}, b_{\mathrm{DNS}}) + \beta \cdot \frac{1}{N}\sum_{i=1}^{N} \sum_{c \in \{11,22,33\}} \left[\mathrm{ReLU}\!\left(-b_{cc}^{(i)} - \tfrac{1}{3}\right)^2 + \mathrm{ReLU}\!\left(b_{cc}^{(i)} - \tfrac{2}{3}\right)^2\right],
    \label{eq:pi_loss}
\end{equation}
which enforces the necessary realizability conditions $-1/3 \leq b_{ii} \leq 2/3$ on the diagonal components.
The penalty weight $\beta$ is linearly ramped from 0 to its target value over the first 100 epochs.
Validation loss is always computed as pure MSE (without penalty terms) for fair comparison.
We sweep $\beta \in \{0.001, 0.01\}$; results are presented in Section~\ref{sec:apriori}.

\subsection{TBRF training}

The TBRF \citep{Kaandorp2020} trains 10 independent random forests (one per basis coefficient) using scikit-learn.
Before training, the 10 tensor basis coefficients $g^{(n)}$ are extracted for each training point via the minimum-norm least-squares solution:
\begin{equation}
    \bm{g} = A^T (A A^T + \epsilon I)^{-1} \bm{b},
\end{equation}
where $A = [T^{(1)} \cdots T^{(10)}]^T$ is the $6 \times 10$ basis matrix and $\epsilon = 10^{-10}$ provides regularization.
Each forest uses 200 trees of maximum depth 20 with default scikit-learn hyperparameters and a fixed random seed of 42.
Compact variants with 1, 5, and 10 trees per coefficient are exported for solver integration experiments.

\subsection{Solver integration}
\label{sec:solver_integration}

The feature computation and inference pipeline within the solver proceeds as follows at each time step:

\begin{enumerate}
    \item \textbf{Velocity gradients}: Compute $\partial U_i / \partial x_j$ (9 components) via second-order central differences on the staggered grid---a standard stencil operation shared with EARSM.
    \item \textbf{Strain/rotation tensors}: Symmetrize and antisymmetrize the gradient tensor, non-dimensionalize by $k/\varepsilon$.
    \item \textbf{Invariants}: Compute the 5 scalar invariants via $3 \times 3$ matrix products and traces.
    \item \textbf{Normalization}: Apply the saved Z-score parameters ($\mu$, $\sigma$ from training).
    \item \textbf{Tensor basis} (TBNN/TBRF only): Compute the 10 basis tensors $T^{(n)}_{ij}$.
    \item \textbf{Inference}: Forward pass through the neural network (MLP/TBNN) or tree traversal (TBRF).
    \item \textbf{Postprocessing}: Write the predicted $\nut$ (MLP) or $\tau_{ij}$ (TBNN/TBRF) to the solver's turbulence fields.
\end{enumerate}

Steps 1--5 are implemented as GPU kernels with each cell processed independently.
For neural network inference (step 6), the computation is parallelized over (cell, neuron) pairs: each GPU thread block processes a batch of cells, with threads within a block computing different output neurons of the same layer.
Weight matrices are stored in constant memory (MLP) or read from global memory with L2 caching (TBNN).
Workspace buffers use a ping-pong pattern between layers to avoid extra memory allocation.

All sub-phases are instrumented with \texttt{TIMED\_SCOPE} profiling markers, enabling the detailed cost breakdown presented in Section~\ref{sec:cost}.

\section{A Priori Evaluation}
\label{sec:apriori}

\subsection{Validation and test set accuracy}

Table~\ref{tab:apriori_rmse} reports the root-mean-square error (RMSE) on the validation and test sets for all models.
The TBNN and TBRF models predict the full $6$-component anisotropy tensor $b_{ij}$, while the MLP models predict a scalar proxy (anisotropy magnitude); RMSE values are therefore not directly comparable across model types but are consistent within each group.

\begin{table}[htbp]
\centering
\caption{A priori accuracy: validation and test set RMSE. TBNN/PI-TBNN/TBRF RMSE is computed over all 6 anisotropy components; MLP RMSE is for the scalar target. Best tensor-predicting model on each metric in bold. The test set comprises periodic hills at $\alpha=1.2$ (interpolation within the training geometry class) and curved backward-facing step at $\Rey=13{,}700$ (unseen geometry).}
\label{tab:apriori_rmse}
\begin{tabular}{lrrrrrr}
\toprule
Model & Val RMSE & Test RMSE & PH $\alpha{=}1.2$ & CBFS & Degradation & Epochs \\
\midrule
TBRF (200 trees) & \textbf{0.0637} & \textbf{0.1253} & \textbf{0.0585} & \textbf{0.1435} & $1.97\times$ & n/a \\
PI-TBNN          & 0.1203          & 0.1418          & 0.0999          & 0.1553          & $1.18\times$ & 729 \\
MLP-Large        & 0.1045          & 0.1772          & 0.0797          & 0.2034          & $1.70\times$ & 344 \\
MLP              & 0.1096          & 0.1843          & 0.0886          & 0.2106          & $1.68\times$ & 1{,}000 \\
TBNN             & 0.0845          & 0.3889          & 0.0755          & 0.4573          & $4.60\times$ & 694 \\
\bottomrule
\end{tabular}
\end{table}

Among tensor-predicting models, the TBRF (200 trees) achieves the lowest validation RMSE (0.064), followed by the TBNN (0.085, 33\% higher).
The 10-tree compact TBRF captures 96\% of the full forest's accuracy improvement over TBNN, suggesting diminishing returns beyond $\sim$10 trees.

The most striking result in Table~\ref{tab:apriori_rmse} is the \emph{reversal} of model rankings between validation and test sets.
The TBNN achieves the best validation RMSE among neural networks (0.085) but the worst test RMSE (0.389)---a $4.6\times$ degradation.
In contrast, the PI-TBNN degrades only $1.18\times$ (0.120 $\to$ 0.142), and the TBRF degrades $1.97\times$ (0.064 $\to$ 0.125).
This reversal is driven almost entirely by the CBFS case (unseen geometry): the TBNN's CBFS RMSE is 0.457, compared to 0.076 on PH $\alpha=1.2$ (an interpolation case within the periodic hills geometry class).

\subsubsection{Per-case breakdown}

Table~\ref{tab:apriori_percase} decomposes the test set RMSE by case.
All models perform well on PH $\alpha=1.2$, which interpolates between training cases at $\alpha=0.5$, $0.8$, $1.0$, and $1.5$.
The CBFS geometry, however, introduces a recirculation region with qualitatively different separation and reattachment physics, exposing each model's ability to extrapolate.

\begin{table}[htbp]
\centering
\caption{Per-case test set RMSE for tensor-predicting models. PH $\alpha=1.2$ is an interpolation case; CBFS $\Rey=13{,}700$ is an unseen geometry. The CBFS/PH ratio quantifies how much worse each model performs on the new geometry.}
\label{tab:apriori_percase}
\begin{tabular}{lrrr}
\toprule
Model & PH $\alpha{=}1.2$ & CBFS $\Rey{=}13{,}700$ & CBFS/PH ratio \\
\midrule
TBRF     & 0.0585 & 0.1435 & $2.45\times$ \\
PI-TBNN  & 0.0999 & 0.1553 & $1.55\times$ \\
TBNN     & 0.0755 & 0.4573 & $6.05\times$ \\
\bottomrule
\end{tabular}
\end{table}

The TBNN's CBFS/PH ratio of $6.05\times$ indicates catastrophic overfitting to the training geometries.
The PI-TBNN, despite sharing the same architecture and parameter count, achieves a CBFS/PH ratio of only $1.55\times$---the most geometry-robust neural network model.
The TBRF's $2.45\times$ ratio reflects the ensemble averaging inherent in random forests, which provides natural regularization against overfitting.

\subsection{Generalization and overfitting}
\label{sec:generalization}

The validation-to-test degradation patterns in Table~\ref{tab:apriori_rmse} reveal a fundamental accuracy--generalization tradeoff.
Figure~\ref{fig:val_test_gap} summarizes this tradeoff: the TBNN sits in the high-accuracy, poor-generalization corner, while the PI-TBNN sacrifices validation accuracy for dramatically better test-set robustness.

The TBNN's overfitting is component-specific.
On the validation set, the TBNN's component-wise RMSE values are uniformly low (e.g., $b_{12}$: 0.053, $b_{11}$: 0.147).
On the CBFS test case, $b_{11}$ RMSE explodes to 0.806 ($5.5\times$ degradation) and $b_{22}$ to 0.580 ($6.8\times$), while $b_{12}$ degrades more modestly to 0.279 ($5.3\times$).
This suggests the TBNN has memorized geometry-specific normal stress magnitudes without learning the underlying physics governing separation and reattachment.

The PI-TBNN's regularization mechanism is noteworthy.
On the validation set, the realizability penalty appears to hurt accuracy ($+42\%$ RMSE vs TBNN).
However, this penalty constrains the network weights to produce outputs near the realizability boundary, limiting the effective capacity of the network and preventing it from fitting geometry-specific artifacts in the training data.
The result is a network that generalizes far better: PI-TBNN's CBFS RMSE (0.155) is $2.9\times$ better than the TBNN's (0.457).
This is a classic bias--variance tradeoff: the PI-TBNN accepts higher bias (worse validation fit) for dramatically lower variance (better generalization).

\subsection{Component-wise accuracy}

Table~\ref{tab:component_rmse} reports the component-wise RMSE on validation and test sets for all tensor-predicting models.

\begin{table}[htbp]
\centering
\caption{Component-wise RMSE on the test set for tensor-predicting models. Values in parentheses are the validation set RMSE for comparison. Components $b_{13}$ and $b_{23}$ are near-zero in the 2D-mean flows and are omitted.}
\label{tab:component_rmse}
\begin{tabular}{lrrr}
\toprule
Component & TBRF & TBNN & PI-TBNN \\
\midrule
$b_{11}$ & 0.167 (0.108) & 0.684 (0.147) & 0.203 (0.220) \\
$b_{12}$ & 0.071 (0.043) & 0.238 (0.053) & 0.074 (0.053) \\
$b_{22}$ & 0.129 (0.066) & 0.492 (0.085) & 0.172 (0.136) \\
$b_{33}$ & 0.211 (0.078) & 0.373 (0.104) & 0.211 (0.128) \\
\bottomrule
\end{tabular}
\end{table}

Several patterns emerge.
First, the shear stress component $b_{12}$ is predicted most accurately by all models, both on validation and test sets, consistent with its direct connection to the mean velocity gradient via the Boussinesq hypothesis.
Second, the TBNN's test-set degradation is concentrated in the normal stress components ($b_{11}$, $b_{22}$, $b_{33}$), where it shows $4$--$7\times$ degradation from validation to test.
Third, the PI-TBNN's $b_{12}$ test RMSE (0.074) is comparable to the TBRF (0.071) and far better than the TBNN (0.238), demonstrating that the realizability penalty preserves shear stress accuracy while preventing normal stress overfitting.

\subsection{Scatter plots}

\todo{Add predicted vs.\ true $b_{ij}$ scatter plots on the test set for key components ($b_{12}$, $b_{11}$) across all tensor-predicting models. Include $R^2$ values.}

\subsection{Lumley triangle}

\todo{Map predicted anisotropy eigenvalues onto the Lumley triangle \citep{Lumley1978} for each model. This provides a visual check of realizability: predictions outside the triangle violate physical constraints on the Reynolds stress tensor.}

\subsection{PI-TBNN beta sweep}
\label{sec:pi_sweep}

The realizability penalty weight $\beta$ in Eq.~\eqref{eq:pi_loss} was swept to determine whether physics-informed constraints improve accuracy.
Table~\ref{tab:pi_sweep} summarizes the results.

\begin{table}[htbp]
\centering
\caption{PI-TBNN realizability penalty sweep. Validation RMSE is pure MSE (no penalty terms).}
\label{tab:pi_sweep}
\begin{tabular}{rrrll}
\toprule
$\beta$ & Val RMSE & Epochs & $\Delta$ vs TBNN & Observation \\
\midrule
0 (TBNN)   & 0.0845 & 694 & ---      & Baseline \\
0.001      & 0.0852 & 729 & $+0.8\%$  & Negligible difference \\
0.01       & 0.0909 & 676 & $+7.6\%$  & Penalty degrades accuracy \\
\bottomrule
\end{tabular}
\end{table}

The unconstrained TBNN already produces near-realizable outputs.
Post-hoc analysis of TBNN predictions on the validation set shows that only 1.29\% of points (309 out of 23{,}967) violate the $b_{ii} \geq -1/3$ bound, and only 0.004\% violate $b_{ii} \leq 2/3$.
The worst violation is $b_{ii} = -0.45$ versus the limit of $-0.33$, and the total realizability penalty magnitude is 0.25\% of the MSE loss.

Based on validation accuracy alone, the penalty appears harmful---consistent with the findings of \citet{Wu2018} that loss-based penalties are redundant given architectural constraints.
However, the test set results (Table~\ref{tab:apriori_rmse}) dramatically reverse this conclusion: the PI-TBNN's test RMSE (0.142) is $2.7\times$ better than the TBNN's (0.389), despite its higher validation RMSE.
The realizability penalty acts as a regularizer that constrains the network's effective capacity, preventing overfitting to geometry-specific features in the training data (see Section~\ref{sec:generalization} for detailed analysis).

\subsection{TBRF tree count sweep}

Table~\ref{tab:tbrf_sweep} shows how TBRF accuracy varies with the number of trees per coefficient.

\begin{table}[htbp]
\centering
\caption{TBRF compact variants: accuracy vs.\ model size.}
\label{tab:tbrf_sweep}
\begin{tabular}{rrrrr}
\toprule
Trees & Total nodes & Binary size & Val RMSE & vs.\ TBNN \\
\midrule
1   & 282{,}902    & \SI{5.7}{MB}   & 0.0778 & $7.9\%$ better \\
5   & 1{,}432{,}612 & \SI{28.7}{MB}  & 0.0678 & $19.8\%$ better \\
10  & 2{,}797{,}130 & \SI{55.9}{MB}  & 0.0650 & $23.1\%$ better \\
200 & 55{,}051{,}026 & \SI{1{,}101}{MB} & 0.0637 & $24.6\%$ better \\
\bottomrule
\end{tabular}
\end{table}

The accuracy improvement saturates rapidly: 10 trees capture 96\% of the full 200-tree ensemble's improvement over TBNN, while requiring only 5\% of the model size.
However, even the 1-tree variant (\SI{5.7}{MB}) involves branching tree traversals that are poorly suited for GPU execution, compared to the TBNN's \SI{196}{KB} of dense matrix multiplies.
This tradeoff between offline accuracy and online deployability is central to the cost analysis in Section~\ref{sec:cost}.

\section{A Posteriori Evaluation}
\label{sec:aposteriori}

To assess whether improved offline accuracy translates to better coupled CFD predictions, we evaluate all models within the solver on three canonical flow configurations.
Each model uses the same spatial discretization, time integrator, and Poisson solver; only the turbulence closure differs.

\subsection{Channel flow at $\Reytau = 180$}

\subsubsection{Setup}

\todo{Specify grid ($128 \times 96 \times 128$), domain size ($4\pi \times 2 \times 2\pi$), $\Rey$ (or $\nu$, $dp/dx$), time integrator (RK3), number of steps, CFL, stretch parameter. Reference: MKM DNS \citep{MKM1999}.}

\subsubsection{Mean velocity profiles}

\todo{Plot $u^+(y^+)$ for all models vs.\ MKM DNS data. Report the $L_2$ error in $u^+$ for each model. Include the law of the wall ($u^+ = y^+$, log law $u^+ = (1/\kappa)\ln y^+ + B$) for reference.}

\subsubsection{Reynolds stress profiles}

\todo{Plot $-\langle u'v' \rangle / u_\tau^2$ vs.\ $y/\delta$ for all models vs.\ MKM DNS. Report component-wise $L_2$ errors. This is the most discriminating metric for tensor-predicting models (TBNN, TBRF) vs.\ scalar models (MLP, SST).}

\subsubsection{Predicted $\Reytau$}

\todo{Report the friction Reynolds number $\Reytau = u_\tau \delta / \nu$ achieved by each model. MKM target: $\Reytau = 178.12$.}

\subsection{Periodic hills at $\Rey_H = 5600$}

\subsubsection{Setup}

\todo{Specify grid, domain, $\Rey$, hill geometry (IBM), reference data \citep{Breuer2009}. Note that periodic hills at different slope angles were in the training set---this tests generalization to a different Reynolds number.}

\subsubsection{Velocity profiles}

\todo{Plot streamwise velocity $u/U_b$ vs.\ $y/H$ at key $x$-stations ($x/H = 0.5, 2, 4, 6, 8$) for all models vs.\ Breuer LES. These stations capture the recirculation zone, reattachment, and recovery regions.}

\subsubsection{Separation and reattachment}

\todo{Report separation and reattachment $x$-locations for each model. Breuer reference: separation at $x/H \approx 0.2$, reattachment at $x/H \approx 4.7$ for $\Rey_H = 5600$.}

\subsubsection{Separation bubble visualization}

\todo{Contour plots or streamlines of the mean flow field showing the separation bubble for selected models (baseline, SST, TBNN) vs.\ reference.}

\subsection{Cylinder at $\Rey = 100$}

\subsubsection{Setup}

\todo{Specify grid ($128 \times 64 \times 64$), domain, cylinder diameter and location (IBM), reference values. Note that the cylinder geometry was not in the McConkey training set---this tests generalization to an entirely unseen geometry.}

\subsubsection{Drag and lift coefficients}

\todo{Report time-averaged drag coefficient $C_d$, RMS lift coefficient $C_{l,\mathrm{rms}}$, and Strouhal number $\mathrm{St}$ for each model. Established reference values: $C_d \approx 1.35$, $C_{l,\mathrm{rms}} \approx 0.34$, $\mathrm{St} \approx 0.165$.}

\subsubsection{Wake profiles}

\todo{Plot streamwise velocity and Reynolds stress profiles in the near wake ($x/D = 1, 2, 5$) for selected models vs.\ reference data.}

\subsection{Grid sensitivity}

\todo{Run the TBNN and SST models on 4 grids of increasing resolution for the channel flow case to demonstrate grid convergence. Report the $L_2$ error in $u^+(y^+)$ vs.\ grid spacing.}

\section{Computational Cost Analysis}
\label{sec:cost}

\subsection{GPU profiling results}

All timing measurements are performed on an NVIDIA L40S GPU with 50{,}000 time steps, excluding the first 100 steps (warm-up).
Table~\ref{tab:gpu_profiling_cyl} reports the results for cylinder flow at $\Rey = 100$ on a $128 \times 64 \times 64$ grid (524K cells), and Table~\ref{tab:gpu_profiling_airfoil} for airfoil flow at $\Rey = 1000$ on a $256 \times 128 \times 128$ grid (4.2M cells).

\begin{table}[htbp]
\centering
\caption{GPU profiling: cylinder $\Rey = 100$, $128 \times 64 \times 64$ (524K cells), L40S, 50K steps.}
\label{tab:gpu_profiling_cyl}
\begin{tabular}{lrrrr}
\toprule
Model & Turb update (s) & Total step (s) & Overhead & Turb \% of total \\
\midrule
Baseline        & 2.33  & 45.0  & ---       & 5.2\% \\
$k$-$\omega$    & 0.47  & 45.1  & $+0.2\%$  & 1.0\% \\
SST             & 1.26  & 46.9  & $+4.2\%$  & 2.7\% \\
EARSM           & 1.46  & 47.1  & $+4.7\%$  & 3.1\% \\
NN-MLP          & 12.78 & 55.6  & $+23.6\%$ & 23.0\% \\
NN-TBNN         & 85.55 & 129.0 & $+186.7\%$ & 66.3\% \\
\bottomrule
\end{tabular}
\end{table}

\begin{table}[htbp]
\centering
\caption{GPU profiling: airfoil $\Rey = 1000$, $256 \times 128 \times 128$ (4.2M cells), L40S, 50K steps.}
\label{tab:gpu_profiling_airfoil}
\begin{tabular}{lrrrr}
\toprule
Model & Turb update (s) & Total step (s) & Overhead & Turb \% of total \\
\midrule
Baseline        & 3.01  & 83.0  & ---       & 3.6\% \\
$k$-$\omega$    & 0.58  & 86.0  & $+3.6\%$  & 0.7\% \\
SST             & 2.14  & 90.0  & $+8.4\%$  & 2.4\% \\
EARSM           & 2.80  & 90.8  & $+9.4\%$  & 3.1\% \\
NN-MLP          & 40.77 & 124.3 & $+49.8\%$ & 32.8\% \\
NN-TBNN         & 543.91 & 628.4 & $+657.1\%$ & 86.6\% \\
\bottomrule
\end{tabular}
\end{table}

Several key observations emerge from these data:

\paragraph{Transport models add minimal overhead.}
The $k$-$\omega$, SST, and EARSM closures add only 0.2--9.4\% to the total step time.
Their stencil-based operations match the GPU's memory access patterns and scale proportionally with the base solver.

\paragraph{NN closures are significantly more expensive.}
The MLP adds 24--50\% overhead, while the TBNN adds 187--657\%.
The TBNN consumes 66--87\% of the total step time, making it the dominant cost rather than the Poisson solve or advection.

\paragraph{NN overhead grows with grid size.}
The TBNN's share of total step time increases from 66\% at 524K cells to 87\% at 4.2M cells.
This occurs because the TBNN's per-cell cost ($\sim$10{,}300~FLOPs) scales linearly with the number of cells, while the base solver's FFT Poisson solve scales as $O(N \log N)$ but with a smaller constant.
The MLP shows a similar trend (23\% to 33\%).
In contrast, transport model overhead remains roughly constant (3--4\%) across grid sizes.

\subsection{Sub-phase timing breakdown}

\todo{Collect and report the per-phase timing breakdown within the turbulence update for each NN model:
\begin{itemize}
    \item Gradient computation (stencil on velocity field)
    \item Feature extraction (invariants from gradients)
    \item Tensor basis computation (10 matrix products)---TBNN/TBRF only
    \item NN forward pass (layer-by-layer matmul + activation)
    \item Postprocessing (write $\nut$ / $\tau_{ij}$)
\end{itemize}
This data is partially available from \texttt{TIMED\_SCOPE} instrumentation (\texttt{nn\_mlp\_features\_gpu}, \texttt{nn\_mlp\_inference\_gpu}, \texttt{nn\_mlp\_postprocess\_gpu}). Need to collect these in formal paper runs.}

\subsection{Why TBNN costs $7.5\times$ more than MLP}

The TBNN's computational cost is dominated by two factors:

\begin{enumerate}
    \item \textbf{Network size}: The TBNN has three $64 \times 64$ hidden layers, each requiring 4{,}096 multiply-add operations per cell, compared to the MLP's single $32 \times 32$ layer (1{,}024 operations). The two large hidden layers alone account for 8{,}192 of the TBNN's 9{,}354 FLOPs (87\%).

    \item \textbf{Tensor basis computation}: The TBNN additionally computes 10 basis tensors $T^{(n)}_{ij}$, requiring $\sim$900~FLOPs of $3 \times 3$ matrix algebra per cell. The MLP does not require this step.

    \item \textbf{Activation functions}: The TBNN evaluates 192 $\tanh$ functions per cell (3 hidden layers $\times$ 64 neurons) versus 64 for the MLP (2 layers $\times$ 32). Each $\tanh$ costs approximately $5\times$ more than a multiply on GPU.
\end{enumerate}

\subsection{Scaling with grid size}

\todo{Run baseline + MLP + TBNN on grids of $64^3$, $128^3$, and $256^3$ cells. Plot turbulence update time vs.\ number of cells. Compare the slope against the theoretical $O(N)$ scaling (linear in cells). Overlay the base solver cost scaling ($O(N)$ for stencils, $O(N \log N)$ for FFT Poisson). This determines whether the NN overhead is a constant fraction of total cost or grows with problem size.}

\subsection{CPU results}

\todo{Run all models on CPU (same grids, same problems) using the GCC build without GPU offload. Report:
\begin{itemize}
    \item CPU wall-clock times for each model
    \item GPU/CPU speedup ratio for the base solver vs.\ for each turbulence model
    \item Whether the cost ranking changes on CPU (e.g., MLP matmuls may be relatively faster on CPU via optimized BLAS, while TBRF tree traversals may suffer less from the absence of warp divergence)
\end{itemize}
}

\subsection{Cost-accuracy tradeoff}
\label{sec:pareto}

\todo{Generate the Pareto plot:
\begin{itemize}
    \item $x$-axis: computational cost (total step time overhead \%, or absolute ms/step)
    \item $y$-axis: a posteriori prediction error (e.g., $L_2$ norm of $u^+(y^+)$ error vs.\ DNS for channel flow)
    \item Each model is a labeled point
    \item Identify the Pareto frontier---models not dominated by any other
    \item Show for multiple flow cases to test robustness of the ranking
    \item Generate separate plots for GPU and CPU
\end{itemize}
This is the central figure of the paper.}

\todo{Generate the cost breakdown stacked bar chart:
\begin{itemize}
    \item For each model, show time spent in each phase (gradient, features, tensor basis, inference, transport, postprocessing)
    \item This visualizes \emph{where} the cost comes from for each model class
\end{itemize}
}

\section{Discussion}
\label{sec:discussion}

\subsection{When to use each model class}

The results in Sections~\ref{sec:apriori}--\ref{sec:cost} suggest a practical decision framework for selecting a turbulence closure:

\begin{itemize}
    \item \textbf{Algebraic / transport models} ($k$-$\omega$ SST, EARSM): For applications where computational budget is the primary constraint---design optimization requiring thousands of evaluations, parametric sweeps, or real-time simulation. These models add less than 10\% overhead on GPU and remain the most robust choice for attached or mildly separated flows.

    \item \textbf{MLP (small)}: A reasonable first step toward data-driven closures when the user wants modest accuracy improvement without dramatic cost increase. At 24--50\% overhead, the MLP is affordable for most applications, though its scalar output limits its ability to capture Reynolds stress anisotropy.

    \item \textbf{PI-TBNN}: The recommended choice when the target geometry may differ from the training set---the common case in engineering practice. Despite 42\% worse validation accuracy than the unconstrained TBNN, the PI-TBNN's test RMSE (0.142) is $2.7\times$ better, with particularly strong robustness on the unseen CBFS geometry (CBFS/PH ratio of $1.55\times$ vs.\ $6.05\times$ for TBNN). The computational cost is identical to the TBNN.

    \item \textbf{TBNN}: Should be used with caution. While it achieves the best validation accuracy among neural networks (RMSE 0.085), it exhibits catastrophic overfitting on unseen geometries (test RMSE 0.389). Only recommended when the deployment geometry closely matches the training data.

    \item \textbf{TBRF}: Not recommended for online use but valuable as an offline reference. The TBRF achieves the best test RMSE (0.125) with good generalization ($1.97\times$ degradation), confirming that the accuracy gap between neural networks and tree ensembles persists on unseen geometries. The ensemble averaging provides natural regularization analogous to the PI-TBNN's penalty.
\end{itemize}

\subsection{PI-TBNN: realizability as regularization}

The realizability penalty in the PI-TBNN tells two very different stories depending on whether one examines validation or test performance (Section~\ref{sec:pi_sweep} and~\ref{sec:generalization}).
On the validation set, the penalty appears harmful: PI-TBNN's RMSE (0.120) is 42\% worse than the unconstrained TBNN (0.085).
On the test set, the conclusion reverses: PI-TBNN (0.142) is $2.7\times$ better than TBNN (0.389).

This dramatic reversal reveals that the realizability penalty acts primarily as a \emph{regularizer}, not as a physics constraint.
The TBNN architecture already embeds significant physical structure:
\begin{itemize}
    \item \textbf{Galilean invariance}: The input features (5 invariants of $\Shat$ and $\Omhat$) are frame-indifferent by construction.
    \item \textbf{Symmetry and zero trace}: The tensor basis contraction $b_{ij} = \sum_n g^{(n)} T^{(n)}_{ij}$ guarantees symmetric, traceless output.
    \item \textbf{Near-realizability}: Only 1.3\% of TBNN validation predictions violate realizability bounds.
\end{itemize}

Given that realizability violations are already rare, the penalty's main effect is to constrain the network's effective capacity.
By penalizing outputs near the boundary of the realizable region, it prevents the coefficient functions $g^{(n)}(\lambda_1,\ldots,\lambda_5)$ from developing sharp, geometry-specific features that fit the training data but do not transfer.
The result is a smoother mapping from invariants to tensor coefficients---higher bias but dramatically lower variance.

This reinterpretation has practical implications.
The earlier conclusion of \citet{Wu2018}---that loss-based physics penalties are redundant given architectural constraints---holds for in-distribution accuracy but is misleading for generalization.
When the deployment domain differs from training (as is typical in engineering practice), the PI-TBNN's regularization effect is valuable precisely because the penalty is ``redundant'' for the training data: it constrains the network without distorting the loss landscape in regions where data is available.

\subsection{TBNN overfitting to training geometries}

The TBNN's catastrophic test-set degradation ($4.6\times$ increase in RMSE) is concentrated on the CBFS geometry, where it achieves an RMSE of 0.457 compared to 0.076 on PH $\alpha=1.2$ (Table~\ref{tab:apriori_percase}).
The PH case interpolates between training geometries ($\alpha=0.5$, $0.8$, $1.0$, $1.5$), while the CBFS introduces a qualitatively different separation mechanism (sharp backward-facing step with streamline curvature).

The component-wise analysis (Table~\ref{tab:component_rmse}) reveals that the overfitting is concentrated in the normal stress components ($b_{11}$, $b_{22}$), which show $5$--$7\times$ degradation from validation to test.
The shear stress $b_{12}$, while also degraded, fares better ($4.5\times$).
This pattern is consistent with the TBNN learning geometry-specific normal stress magnitudes---which vary strongly between flow configurations---rather than the universal relationship between strain/rotation invariants and stress anisotropy.

The TBNN has only 9,354 parameters (identical to the PI-TBNN), so the overfitting is not driven by excessive model capacity in the traditional sense.
Rather, the unconstrained loss function allows the network to exploit correlations between invariant features and geometry-specific stress magnitudes that happen to be present in the training set but do not generalize.
The PI-TBNN's realizability penalty disrupts exactly these spurious correlations.

\subsection{TBRF: offline ceiling, online impractical}

The TBRF achieves the highest offline accuracy (RMSE = 0.064 vs.\ 0.085 for TBNN, a 24.6\% improvement) but represents a fundamentally different computational paradigm:
\begin{itemize}
    \item Decision tree traversals involve \emph{data-dependent branching}: different cells follow different paths through the tree, causing warp divergence on GPU.
    \item Tree node data (feature indices, thresholds, child pointers) is scattered in memory with no spatial locality, causing cache misses.
    \item The 10-tree model requires \SI{56}{MB} of tree data, exceeding GPU L2 cache capacity.
\end{itemize}

These are not implementation limitations that can be overcome with better engineering; they reflect a fundamental mismatch between the random forest algorithm and GPU hardware.
The TBRF's value extends beyond establishing a validation accuracy ceiling.
On the test set, the TBRF achieves the best RMSE (0.125) among all models, with a val-to-test degradation ratio of $1.97\times$---far better than the TBNN ($4.60\times$) though not as robust as the PI-TBNN ($1.18\times$).
The ensemble averaging inherent in random forests provides natural regularization, analogous to the PI-TBNN's realizability penalty but through a different mechanism (model averaging vs.\ capacity constraint).
This suggests that the TBRF's offline accuracy advantage over neural networks is partly an accuracy advantage and partly a generalization advantage---it fits the training data better \emph{and} transfers better to new geometries.

\subsection{CPU vs.\ GPU cost differences}

\todo{Discuss how the cost ranking changes between GPU and CPU:
\begin{itemize}
    \item On CPU, matrix multiplies benefit from optimized BLAS libraries, potentially making MLP/TBNN relatively cheaper.
    \item On CPU, tree traversals do not suffer from warp divergence, potentially making TBRF relatively less expensive.
    \item The GPU/CPU speedup ratio may differ across model types: stencil-based models may show higher speedup than NN models.
\end{itemize}
}

\subsection{Limitations}

Several limitations of this study should be noted:

\begin{itemize}
    \item \textbf{Training data dependence}: All NN models are trained on $k$-$\omega$ SST fields from the McConkey dataset. Performance may differ with other RANS baselines ($k$-$\varepsilon$, Spalart--Allmaras) or different training data sources.

    \item \textbf{Generalization}: The test set includes only two cases (periodic hills at $\alpha=1.2$ and CBFS at $\Rey=13{,}700$). While these reveal significant overfitting (Section~\ref{sec:generalization}), performance on dramatically different geometries (3D bluff bodies, internal flows, free shear layers) is not evaluated. The TBNN's catastrophic CBFS failure may or may not be representative of its behavior on other unseen geometries.

    \item \textbf{Single solver}: All cost measurements are from one solver implementation. Different solvers (spectral element, finite volume, unstructured) may show different relative overheads due to different baseline costs.

    \item \textbf{Inference optimization}: Our NN inference uses explicit matrix-vector products. Optimized inference engines (TensorRT, ONNX Runtime) could potentially reduce NN overhead, though the qualitative ordering is unlikely to change.

    \item \textbf{No online training}: We evaluate only fixed-weight, offline-trained models. Online or adaptive training approaches may offer different cost-accuracy tradeoffs.

    \item \textbf{Single GPU type}: Primary results are on the L40S. The relative cost may differ on other GPU architectures (different L2 cache sizes, different FLOP/bandwidth ratios). \todo{Add results on H100/H200 in supplement.}
\end{itemize}

\subsection{Comparison to prior work}

\todo{Discuss how our a priori accuracy compares to reported values in \citet{Ling2016} and \citet{Kaandorp2020}. Note differences in dataset, split protocol, and architecture. Discuss whether the cost ranking changes in the a posteriori setting vs.\ a priori accuracy ordering.}

\section{Conclusions}
\label{sec:conclusions}

\todo{Revise after all results are available.}

We have presented a systematic comparison of five data-driven turbulence closure architectures and four classical RANS models, all implemented in a single GPU-accelerated incompressible Navier--Stokes solver and evaluated on the same training data, test cases, and hardware.
The key findings are:

\begin{enumerate}
    \item \textbf{Classical transport models remain remarkably cheap.}
    The $k$-$\omega$ SST and EARSM closures add only 0.2--9.4\% to the total solver step time on GPU, even for anisotropic Reynolds stress prediction.
    Their stencil-based operations naturally match GPU memory access patterns.

    \item \textbf{Neural network closures carry significant overhead.}
    The MLP (1{,}249 parameters) adds 24--50\% overhead; the TBNN (9{,}354 parameters) adds 187--657\%.
    The overhead grows with grid size because NN inference scales as $O(N)$ with a large per-cell constant, while the base solver benefits from sublinear FFT scaling.
    \todo{Add quantitative a posteriori finding here.}

    \item \textbf{Physics-informed realizability penalties do not improve TBNN accuracy.}
    The tensor basis architecture already constrains outputs to be near-realizable (1.3\% violation rate), making soft penalty terms redundant or harmful ($+7.6\%$ RMSE at $\beta = 0.01$).

    \item \textbf{The TBRF is the most accurate model offline but impractical online.}
    Its 24.6\% RMSE improvement over TBNN comes at the cost of branch-heavy tree traversals that are fundamentally poorly suited for GPU execution.
    The TBRF's primary value is as an accuracy ceiling for the invariant-basis feature representation.
    \todo{Add TBRF solver timing.}

    \item \textbf{The cost--accuracy Pareto frontier identifies \todo{which model} as the recommended choice for GPU-accelerated simulations.}
    \todo{Fill in based on Pareto analysis.}
\end{enumerate}

\subsection{Practical recommendations}

Based on these findings, we offer the following guidance for practitioners:

\begin{itemize}
    \item For \textbf{design optimization} or \textbf{parametric sweeps} (many solver evaluations): use $k$-$\omega$ SST or EARSM. The sub-10\% overhead allows hundreds of evaluations in the time a single TBNN run would take.

    \item For \textbf{high-accuracy predictions of separated or anisotropic flows} (single or few evaluations): the TBNN provides the best accuracy among deployable models. \todo{Quantify accuracy gain vs.\ SST on periodic hills.}

    \item For \textbf{initial exploration} of data-driven closures: the small MLP is an affordable entry point at 24--50\% overhead, suitable for assessing whether a data-driven approach improves results for a given flow class.

    \item \textbf{Avoid deploying TBRF in online solvers}. Use it for offline analysis (accuracy ceiling, feature importance) and consider distilling its predictions into a more compact neural network.
\end{itemize}

\subsection{Future work}

Several directions merit further investigation:

\begin{itemize}
    \item \textbf{Inference optimization}: TensorRT, ONNX Runtime, or custom CUDA kernels could reduce NN overhead by fusing operations and leveraging mixed-precision arithmetic (FP16/BF16 for weights).

    \item \textbf{Architecture search}: Smaller TBNN variants (fewer layers, fewer neurons) or hybrid architectures that predict a low-rank approximation to $b_{ij}$ may offer better cost-accuracy tradeoffs.

    \item \textbf{Knowledge distillation}: Train a compact MLP to approximate TBRF predictions, combining the TBRF's offline accuracy with the MLP's GPU efficiency.

    \item \textbf{Online training}: Adapt model weights during the simulation using solver-generated data, potentially improving generalization without offline training datasets.

    \item \textbf{Multi-GPU scaling}: Investigate how NN inference overhead interacts with MPI domain decomposition and multi-GPU communication.
\end{itemize}


\appendix
\section{Training Methodology}
\label{sec:appendix_training}

This appendix provides the complete training specification for reproducibility.

\subsection{Common hyperparameters}

Table~\ref{tab:hyperparams} lists all hyperparameters shared across neural network models.

\begin{table}[htbp]
\centering
\caption{Training hyperparameters (all neural networks).}
\label{tab:hyperparams}
\begin{tabular}{ll}
\toprule
Hyperparameter & Value \\
\midrule
Optimizer                & Adam \citep{Kingma2015} \\
Initial learning rate    & $10^{-3}$ \\
Adam $\beta_1, \beta_2$  & 0.9, 0.999 \\
Adam $\epsilon$          & $10^{-8}$ \\
LR schedule              & Cosine annealing with warm restarts \\
Cosine $T_0$             & 200 epochs \\
Cosine $T_{\mathrm{mult}}$ & 2 \\
Cosine $\eta_{\min}$     & $10^{-6}$ \\
Batch size               & 256 \\
Max epochs               & 1{,}000 \\
Early stopping patience  & 150 epochs \\
Early stopping criterion & Validation MSE \\
Best model selection     & Lowest validation MSE across all epochs \\
Loss function            & MSE \\
Regularization           & None \\
Data precision (training)    & float32 \\
Data precision (features)    & float64 \\
\bottomrule
\end{tabular}
\end{table}

\subsection{Model architectures}

\begin{table}[htbp]
\centering
\caption{Neural network architecture details.}
\label{tab:architectures}
\begin{tabular}{lllrrl}
\toprule
Model & Layers & Activation & Parameters & Output & Target \\
\midrule
MLP       & $5 \to 32 \to 32 \to 1$         & Tanh / Linear & 1{,}249  & Scalar   & $|b|$ \\
MLP-Large & $5 \to 128 \to 128 \to 128 \to 128 \to 1$ & Tanh / Linear & 50{,}049 & Scalar   & $|b|$ \\
TBNN      & $5 \to 64 \to 64 \to 64 \to 10$ & Tanh / Linear & 9{,}354  & 10 coeff.\ & $b_{ij}$ \\
PI-TBNN   & Same as TBNN                     & Tanh / Linear & 9{,}354  & 10 coeff.\ & $b_{ij}$ \\
\bottomrule
\end{tabular}
\end{table}

\subsection{TBRF hyperparameters}

\begin{table}[htbp]
\centering
\caption{TBRF (scikit-learn \texttt{RandomForestRegressor}) hyperparameters.}
\label{tab:tbrf_hyperparams}
\begin{tabular}{ll}
\toprule
Parameter & Value \\
\midrule
Trees per coefficient    & 200 \\
Number of coefficients   & 10 (Pope basis) \\
Total trees              & 2{,}000 \\
Max depth                & 20 \\
Min samples split        & 2 (default) \\
Min samples leaf         & 1 (default) \\
Max features             & $\sqrt{n_{\mathrm{features}}}$ (default) \\
Bootstrap                & True \\
Random state             & 42 \\
\bottomrule
\end{tabular}
\end{table}

\subsection{Training results}

\begin{table}[htbp]
\centering
\caption{Training outcomes for all models.}
\label{tab:training_results}
\begin{tabular}{lrrrl}
\toprule
Model & Val RMSE & Epochs run & Stopped by & Training time \\
\midrule
TBRF (200 trees)         & 0.0637 & n/a   & n/a         & $\sim$500 s (4 CPU cores) \\
TBNN                     & 0.0845 & 694   & Early stop  & \todo{TBD} \\
PI-TBNN ($\beta=0.001$)  & 0.0852 & 729   & Early stop  & \todo{TBD} \\
PI-TBNN ($\beta=0.01$)   & 0.0909 & 676   & Early stop  & \todo{TBD} \\
MLP-Large                & 0.1045 & 344   & Early stop  & \todo{TBD} \\
MLP                      & 0.1096 & 1{,}000 & Max epochs & \todo{TBD} \\
\bottomrule
\end{tabular}
\end{table}

\subsection{Software and hardware}

\begin{table}[htbp]
\centering
\caption{Software versions used for training.}
\label{tab:software}
\begin{tabular}{ll}
\toprule
Package & Version \\
\midrule
Python     & 3.x (NVIDIA PyTorch container) \\
PyTorch    & 2.1.0a0+32f93b1 \\
NumPy      & 1.22.2 \\
scikit-learn & 1.2.0 \\
pandas     & 1.5.3 \\
\bottomrule
\end{tabular}
\end{table}

Training was performed on an NVIDIA L40S GPU at the Georgia Tech PACE cluster.
No fixed random seed was used for PyTorch (weight initialization and data shuffling are non-deterministic).
Feature extraction used float64 precision; training used float32.

\subsection{Weight export format}

Neural network weights are exported as human-readable text files with 10-digit scientific notation:
\begin{verbatim}
data/models/<model>_paper/
  layer0_W.txt    # [out_features, in_features]
  layer0_b.txt    # [out_features]
  layer1_W.txt
  ...
  input_means.txt # Z-score means [5]
  input_stds.txt  # Z-score stds [5]
  metadata.json   # Architecture info
\end{verbatim}

TBRF compact variants are exported as flat binary files with a 12-byte header (\texttt{total\_nodes}, \texttt{n\_basis}, \texttt{n\_trees} as int32) followed by five arrays (\texttt{children\_left}, \texttt{children\_right}, \texttt{feature}, \texttt{threshold}, \texttt{value}).

\section{Tensor Basis Functions}
\label{sec:appendix_basis}

The 10-tensor Pope integrity basis \citep{Pope1975} is computed from the non-dimensionalized strain rate $\Shat$ and rotation rate $\Omhat$:
\begin{align}
    T^{(1)} &= \Shat \\
    T^{(2)} &= \Shat\Omhat - \Omhat\Shat \\
    T^{(3)} &= \Shat^2 - \tfrac{1}{3}\mathrm{tr}(\Shat^2)I \\
    T^{(4)} &= \Omhat^2 - \tfrac{1}{3}\mathrm{tr}(\Omhat^2)I \\
    T^{(5)} &= \Omhat\Shat^2 - \Shat^2\Omhat \\
    T^{(6)} &= \Omhat^2\Shat + \Shat\Omhat^2 - \tfrac{2}{3}\mathrm{tr}(\Shat\Omhat^2)I \\
    T^{(7)} &= \Omhat\Shat\Omhat^2 - \Omhat^2\Shat\Omhat \\
    T^{(8)} &= \Shat\Omhat\Shat^2 - \Shat^2\Omhat\Shat \\
    T^{(9)} &= \Omhat^2\Shat^2 + \Shat^2\Omhat^2 - \tfrac{2}{3}\mathrm{tr}(\Shat^2\Omhat^2)I \\
    T^{(10)} &= \Omhat\Shat^2\Omhat^2 - \Omhat^2\Shat^2\Omhat
\end{align}
Each tensor is symmetric and traceless.
In the solver, only the 6 independent components $(11, 12, 13, 22, 23, 33)$ are stored and computed.

\section{Complete Profiling Tables}
\label{sec:appendix_profiling}

\todo{Add complete profiling tables for all GPU types (L40S, H100, H200), all flow cases, all models, including sub-phase breakdown.}


\bibliographystyle{plainnat}
\bibliography{references}

\end{document}
